We envision several promising directions to extend this neuro-symbolic architecture:

\textbf{Cross-Modal Knowledge Expansion}: We plan to integrate disparate medical modalities (e.g., CT scans, electronic health records) into the knowledge graph, evolving the benchmark from visual perception to comprehensive clinical diagnosis.
% 跨模态知识扩展：我们计划将不同的医疗模式（例如 CT 扫描、电子健康记录）整合到知识图谱中，将基准从视觉感知发展到全面的临床诊断。

\textbf{Active Learning Feedback Loop}: We aim to close the loop by using the performance metrics of evaluated models to automatically identify "weak paths" in the knowledge graph, guiding the pipeline to generate harder, adversarial training samples for continuous model improvement.
% \textbf{主动学习反馈循环}：我们旨在通过使用评估模型的性能指标来自动识别知识图谱中的“薄弱路径”，从而形成闭环，引导管道生成更难的对抗性训练样本，以持续改进模型。

\textbf{Privacy-Preserving Federated Synthesis}: Given the sensitivity of medical data, we intend to explore federated execution of our pipeline, allowing hospitals to generate standardized benchmarks locally without sharing raw patient video data.
% 隐私保护型联邦合成：鉴于医疗数据的敏感性，我们打算探索我们管道的联邦执行，允许医院在本地生成标准化基准，而无需共享原始患者视频数据。